\begin{lemma}
For the two-bite value \(\beta=1/2\), the untruncated edge probability
\[
p(n)=\noderef{TwoBiteEdgeProbability}(1/2,n)
=\frac12\sqrt{\frac{\log n}{n}}
\]
eventually lies in the interval \([0,1]\).
\end{lemma}
\begin{proof}
For all \(n\ge1\), the real number \(n\) is positive.  For all sufficiently
large \(n\), say \(n\ge3\), we also have \(\log n\ge0\).  Hence
\(\log n/n\ge0\), its square root is nonnegative, and therefore
\[
0\le \frac12\sqrt{\frac{\log n}{n}}.
\]

It remains to prove the upper bound eventually.  The standard logarithmic
comparison \(\log n/n\to0\) gives, for all sufficiently large \(n\),
\[
\frac{\log n}{n}\le4.
\]
Taking square roots preserves the inequality because both sides are
nonnegative, so
\[
\sqrt{\frac{\log n}{n}}\le2.
\]
Multiplying by \(1/2\) gives \(p(n)\le1\).  Combining the eventual lower and
upper bounds proves the stated eventual membership in \([0,1]\).
\end{proof}
